% !Mode:: "TeX:UTF-8"
% 任务书中的信息
%% 原始资料及设计要求
\assignReq
{1．查阅深度学习框架测试、模糊测试、向量检索等相关中英文文献。}
{2．以PyTorch模糊测试场景为背景，分析测试用例、日志与覆盖率数据特点。}
{3．设计测试用例管理与分析系统，完成用例存储、检索与分析。}
{4．系统采用前后端分离架构，支持候选边界条件分析与结果展示。}
{5．论文撰写符合本科毕业设计规范，提交软件、论文及相关材料。}
%% 工作内容
\assignWork
{1．完成课题调研与需求分析，明确系统目标与技术路线。}
{2．完成测试数据处理流程设计，实现日志解析与特征提取。}
{3．实现用例管理与分析功能，包括语义检索、故障辅助定位与候选边界条件分析。}
{4．完成系统前后端开发与集成，实现列表展示、详情分析与可视化。}
{5．搭建实验环境，基于真实测试数据开展实验并分析结果。}
{6．完成论文撰写、修改定稿及答辩准备工作。}
%% 参考文献
\assignRef
{[1] Pei K, Cao Y, Yang J, et al. DeepXplore:}
{\hspace*{1.8em}Automated Whitebox Testing of Deep Learning Systems[C]. SOSP, 2017.}
{[2] Xie X, Ma L, Juefei-Xu F, et al. DeepHunter:}
{\hspace*{1.8em}A Coverage-Guided Fuzz Testing Framework for DNNs[C]. ISSTA, 2019.}
{[3] Ernst M D, Perkins J H, Guo P J, et al. The Daikon System:}
{\hspace*{1.8em}Dynamic Detection of Likely Invariants[J]. Sci. Comput. Program., 2007.}
{[4] Lewis P, Perez E, Piktus A, et al. Retrieval-Augmented Generation}
{\hspace*{1.8em}for Knowledge-Intensive NLP Tasks[C]. NeurIPS, 2020.}
